% Supplementary material for the computer-assisted accuracy proof.
%
% This is a new additive file.  It consolidates the implemented Julia/Arb
% pipeline, the Path-A/Path-B Hilbert-transform architecture, and the audited
% Lean interface for the phase-4 N=150 half-DC certificate.  The 2026-09-09
% revision adds paper-style theorem packaging, implementation-entry tables,
% and Crunchy2 run metadata inspired by the level of detail in
% Buckmaster--Cao-Labora--Gomez-Serrano's computer-assisted appendix.
% It does not replace the earlier interval-arithmetic or Lean supplementary
% drafts.
%
% Principal repository sources:
%   docs/overview.md
%   docs/julia_design.md (architecture only; implementation details are stale)
%   docs/lean_design.md
%   docs/lean_roadmap_2026-05-23.md
%   docs/sharding_plan_2026-07-03.md
%   docs/half_dc_alignment.md
%   docs/pinn_parameter_convention.md
%   docs/lean_certificate_audit_20260831.md
%   BabenkoRigorous/src/*.jl
%   WaveCAP/WaveCAP/*.lean

\documentclass[11pt]{article}

\usepackage[T1]{fontenc}
\usepackage{amsmath,amssymb}
\usepackage{amsthm}
\usepackage{booktabs}
\usepackage{longtable}
\usepackage{array}
\usepackage[margin=1in]{geometry}
\usepackage{microtype}
\usepackage[hidelinks]{hyperref}
\urlstyle{same}

\title{Supplementary Material:\\Computer-Assisted Accuracy Proof}
\author{}
\date{Version of 2026-10-01}

\newtheorem{theorem}{Supplementary Theorem}
\newtheorem{proposition}[theorem]{Supplementary Proposition}
\newtheorem{lemma}[theorem]{Supplementary Lemma}

\providecommand{\CandidateDegree}{150}
\providecommand{\CandidateCount}{151}
\providecommand{\VerificationGrid}{2{,}000{,}000}
\providecommand{\TaylorOrder}{6}
\providecommand{\WorkingPrecision}{128}
\providecommand{\WaveSpeed}{1.1687864806007486}
\providecommand{\BabenkoMu}{1.3660618372350841}
\providecommand{\PinnBeta}{0.498}
\providecommand{\BabenkoBeta}{0.248004}
\providecommand{\ResidualMidpoint}{1.9928311263803702\times10^{-14}}
\providecommand{\ResidualRadius}{1.5799937398052128\times10^{-13}}
\providecommand{\ResidualLower}{-1.3807106271671758\times10^{-13}}
\providecommand{\ResidualUpper}{1.7792768524432498\times10^{-13}}
\providecommand{\DLower}{3.4465883518530394\times10^{-1}}
\providecommand{\RunArchiveSHA}{4d65357a64a0ccfbf0ed67f8dd0e1970f0ceb549a8cbf647d8490e4ac48abb07}
\providecommand{\HalfDCArchiveSHA}{ee8fe5db1ea3bbd6e37d0888fdadc00c9c9e1d6ce81cb9d46ae3e4db8ffcdf43}
\providecommand{\HalfDCCertificateSHA}{2190a52f3520ab0b15683eab592ee1105797d737a22f2e24a869130f5a431ed2}

\begin{document}
\maketitle

\section{Computer-assisted accuracy proof}
\label{sec:supp-computer-assisted-accuracy}
\begingroup
\setlength{\emergencystretch}{1em}

\subsection{How to read this supplement}
\label{sec:supp-how-to-read}

This supplement is written as a self-contained computer-assisted-proof
appendix.  Each numerical claim is tied to a mathematical statement, the
Julia routine that produced the enclosure, the Lean module that checks the
corresponding logical consequence, and the archived Crunchy2 artifact from
which the final constants were taken.  The organization is intentionally
more explicit than a short Methods section: the goal is to make the boundary
between discovery, interval computation, formal proof, and remaining
assumptions auditable by a reader.

The main theorem exported by the supplement is a residual certificate for a
single finite cosine series.  The theorem is not a nearby-solution theorem,
and no step in this document asserts an inverse estimate for the linearized
Babenko operator.  The supplement therefore records exactly what has been
certified and exactly what would be needed for a later existence upgrade.

\subsection{Statement and scope of the certificate}
\label{sec:supp-cap-scope}

The purpose of the computer-assisted calculation is to certify the accuracy
of one explicit numerical traveling-wave profile.  After the search and
polishing stages, the profile is frozen as the finite cosine polynomial
\begin{equation}
 y_N(u)=\frac{a_0}{2}+\sum_{k=1}^{N}a_k\cos(ku),
 \qquad N=\CandidateDegree.
 \label{eq:supp-cap-candidate}
\end{equation}
The coefficients in Eq.~\eqref{eq:supp-cap-candidate} are fixed binary
floating-point numbers.  The verification treats their exact dyadic values
as the definition of the candidate; it does not assume that they approximate
unknown ideal coefficients.  The neural network is therefore a search tool,
not part of the object appearing in the theorem.

Let \(H\) denote the \(2\pi\)-periodic Hilbert transform with convention
\begin{equation}
 H[\cos(ku)]=\sin(ku),\qquad
 H[\sin(ku)]=-\cos(ku),\qquad H[1]=0.
 \label{eq:supp-cap-H-modes}
\end{equation}
For
\begin{equation}
 D=\sqrt{y_u^2+\bigl(1+H[y_u]\bigr)^2},
 \qquad
 G=\frac{1+H[y_u]}{D},
 \label{eq:supp-cap-DG}
\end{equation}
the residual evaluated by the verifier is
\begin{equation}
\begin{split}
 R(y)={}&\mu H[y_u]-y-yH[y_u]-H[yy_u]\\
 &+\beta\,\partial_u\!\left(\frac{y_u}{D}\right)
 -\beta H[\partial_uG].
\end{split}
\label{eq:supp-cap-residual}
\end{equation}
The PINN uses the parameters \(c\) and \(b\), whereas
Eq.~\eqref{eq:supp-cap-residual} uses
\begin{align}
 \mu&=\operatorname{fl}_{64}(c^2)=\BabenkoMu,\\
 \beta&=\operatorname{fl}_{64}(b^2)=\BabenkoBeta,\\
 c&=\WaveSpeed,\qquad b=\PinnBeta.
 \label{eq:supp-cap-parameters}
\end{align}
Here \(\operatorname{fl}_{64}\) denotes the stored IEEE double-precision
result of the indicated operation.  The certificate later records those
stored values as exact dyadic rationals.  This squaring is part of the
certificate specification.  Substituting the unsquared PINN parameters into
Eq.~\eqref{eq:supp-cap-residual} would verify a different equation.

The calculation proves two continuum inequalities:
\begin{equation}
 \inf_{u\in\mathbb R}D(u)>0,
 \qquad
 \sup_{u\in\mathbb R}|R(y_N)(u)|\leq\varepsilon_*.
 \label{eq:supp-cap-objectives}
\end{equation}
For the successful run,
\begin{equation}
 D(u)\geq0.3446588,
 \qquad
 |R(y_{150})(u)|\leq1.78\times10^{-13}
 \quad\text{for every }u\in\mathbb R.
 \label{eq:supp-cap-readable-result}
\end{equation}
The manuscript Babenko operator \(\mathcal B\) uses the opposite residual
sign in the normalization \(g=1\):
\begin{equation}
 R(y_N)(u)=-\mathcal B(y_N)(u).
 \label{eq:supp-cap-sign}
\end{equation}
Consequently Eq.~\eqref{eq:supp-cap-readable-result} is also the claimed
bound for \(|\mathcal B(y_{150})|\).

This is an accuracy theorem for a supplied finite series.  It is not an
existence or uniqueness theorem for an exact nearby water wave.  Such a
result would additionally require a rigorous inverse bound for the
linearized Babenko operator and a Newton--Kantorovich or contraction
argument; neither ingredient is part of the present certificate.

\subsection{Certified statement and proof entry points}
\label{sec:supp-certified-statement}

\begin{theorem}[Crunchy2 phase--4 \(N=150\) residual certificate]
\label{thm:supp-crunchy2-residual}
Let \(y_{150}\) be the finite cosine polynomial in
Eq.~\eqref{eq:supp-cap-candidate}, with the \CandidateCount{} coefficients
stored in
\texttt{CertificatePhase4ScalarN150}\allowbreak
\texttt{Phase420260827HalfDC.lean}.  Let
\(\mu=\BabenkoMu\) and \(\beta=\BabenkoBeta\) be the exact dyadic values
recorded in the same certificate.  Then the denominator \(D\) from
Eq.~\eqref{eq:supp-cap-DG} is positive for every \(u\in\mathbb R\), and
\begin{equation}
 \forall u\in\mathbb R,\qquad
 |R(y_{150})(u)|\leq
 \frac{178}{10^{15}}.
 \label{eq:supp-crunchy2-theorem}
\end{equation}
Equivalently, for the manuscript Babenko operator \(\mathcal B\) in the
normalization \(g=1\),
\begin{equation}
 \forall u\in\mathbb R,\qquad
 |\mathcal B(y_{150})(u)|\leq
 \frac{178}{10^{15}}.
 \label{eq:supp-crunchy2-paper-theorem}
\end{equation}
\end{theorem}

\begin{proof}[Proof architecture]
The profile and parameters are fixed as exact rational data in the generated
Lean certificate.  Julia/Arb proves the denominator and residual anchors
\(\mathsf A_D\) and \(\mathsf A_R\) described in
Eqs.~\eqref{eq:supp-cap-anchor-D}--\eqref{eq:supp-cap-anchor-R}.  Lean then
checks the denominator corollary, the residual theorem on one period, the
periodic extension to all real \(u\), the exact rational comparison
\(\varepsilon_{0,\rm cert}<178/10^{15}\), and the sign bridge between
\(R\) and \(\mathcal B\).  The two Arb anchors are the only
certificate-specific numerical assumptions in the final theorem.
\end{proof}

\begin{proposition}[Interval computation recorded by the archive]
\label{prop:supp-interval-archive}
The archived Crunchy2 run for
\texttt{phase4\_20260827\_scalar\_N150} produced the Arb ball
\begin{equation}
 B_R=[\ResidualMidpoint\mathbin{\pm}\ResidualRadius],
 \label{eq:supp-recorded-ball}
\end{equation}
whose outward endpoints are
\[
 \ResidualLower\leq B_R\leq\ResidualUpper.
\]
The upper endpoint is the value serialized as
\(\varepsilon_{0,\rm cert}\).  The same run produced the denominator lower
endpoint \(D_{\rm lower}=\DLower\), serialized as
\(D_{\rm lower,cert}\).
\end{proposition}

\begin{proposition}[Formal trust boundary]
\label{prop:supp-formal-trust-boundary}
The exact residual and paper-operator Lean theorems depend on
\texttt{cert\_arb\_dmin\_bound} and
\texttt{cert\_arb\_sup\_bound\_le\_eps0}, together with the standard Mathlib
principles \texttt{propext}, \texttt{Classical.choice}, and
\texttt{Quot.sound}.  The denominator theorem uses only the denominator
anchor and the standard principles.  The half-DC convention lemmas and
paper-sign alignment theorem use no certificate-specific numerical axioms.
\end{proposition}

\begin{center}
\small
\begin{tabular}{@{}>{\raggedright\arraybackslash}p{0.23\linewidth}>{\raggedright\arraybackslash}p{0.31\linewidth}>{\raggedright\arraybackslash}p{0.36\linewidth}@{}}
\toprule
Certified item & Julia or archive entry & Lean entry \\
\midrule
Finite candidate and dyadic parameters &
\path{pinn_a_coeffs_N150_phase4_20260827.txt};
\path{Certificate.jl} &
\texttt{CertificatePhase4ScalarN150}\allowbreak
\texttt{Phase420260827HalfDC.lean} \\
\addlinespace
Path-A Hilbert transform on finite series &
\path{src/FourierSeries.jl} &
\path{Definitions.lean};
\path{HilbertOnTrigPoly.lean} \\
\addlinespace
Path-B sampled Hilbert transform and aliasing &
\path{src/FFTHilbert.jl} &
\path{AliasingErrorBound.lean};
\path{FFTHilbertCorrectness.lean} \\
\addlinespace
Taylor data, products, reciprocal, square root &
\path{src/TaylorSeries.jl};
\path{src/GDerivatives.jl} &
\path{LeibnizBound.lean};
\path{FaaDiBrunoBound.lean};
\path{TaylorSeriesCorrectness.lean} \\
\addlinespace
Denominator positivity &
\texttt{dmin\_bound};
\path{staged_dlower.arbdump} &
\path{DminBoundCorrectness.lean};
\texttt{cert\_arb\_dmin\_bound} \\
\addlinespace
Residual assembly and continuum scan &
\path{src/ResidualBound.jl};
\path{staged_bound.arbdump} &
\path{GridToSupNormTaylor.lean};
\path{ResidualAssemblyCorrectness.lean};
\texttt{cert\_arb\_sup\_bound\_le\_eps0} \\
\addlinespace
Paper theorem and audit &
\path{phase4_scalar_n150_20260827_2m_results_20260829.tgz} &
\texttt{MainPhase4ScalarN150}\allowbreak
\texttt{Phase420260827HalfDC.lean};
\path{PaperAudit.lean} \\
\bottomrule
\end{tabular}
\end{center}

\subsection{Overview of the proof architecture}
\label{sec:supp-cap-overview}

The proof has three conceptual layers.  First, finite Fourier algebra
constructs the candidate and handles the Hilbert transform whenever its
input is a known trigonometric polynomial.  Second, Julia and Arb compute
outward-rounded interval enclosures for the nonlinear residual, add analytic
corrections for unresolved Fourier modes, and use Taylor's theorem to pass
from a grid to the whole torus.  Third, Lean defines the same mathematical
objects, proves the reusable analytic lemmas underlying the calculation, and
checks the exact logical consequences of the exported numerical
certificate.

The Hilbert-transform work is divided into two pathways.
\begin{description}
 \item[Path A: finite trigonometric algebra.]
 The input has a finite Fourier representation, so \(H\) acts
 mode-by-mode using Eq.~\eqref{eq:supp-cap-H-modes}.  No sampling or aliasing
 correction is needed for this operation.

 \item[Path B: certified sampled transform.]
 The input is represented by interval-enclosed values and derivatives on a
 uniform grid, rather than by a finite exact coefficient vector.  A discrete
 Hilbert transform is evaluated by certified convolution, and a rigorous
 Fourier-tail estimate is added to enclose the continuous transform.
\end{description}
Path B is essential for \(H[\partial_uG]\), because the square root and
reciprocal in Eq.~\eqref{eq:supp-cap-DG} generally give \(G\) an infinite
Fourier series.  The current Julia implementation also routes
\(H[yy_u]\) through Path B for uniformity, although \(yy_u\) is itself a
finite trigonometric polynomial.  For the present candidate its degree is at
most 300, far below the Nyquist frequency of the two-million-point grid, so
an exact coefficient treatment would have no unresolved modes.  The generic
Path-B routine does not exploit that special fact and adds its valid but
conservative tail correction.  This implementation choice does not project
the product back to degree \(N\).

\subsection{The finite candidate and nonlinear mode generation}
\label{sec:supp-cap-candidate-construction}

The final-stage PINN fields were sampled on a Fourier grid and projected onto
finite trigonometric series.  Their coefficients initialized Newton's method
for a degree--160 Galerkin system.  After polishing, the auxiliary fields
were discarded and the scalar profile was truncated to degree
\(N=\CandidateDegree\), leaving \CandidateCount{} stored coefficients in
Eq.~\eqref{eq:supp-cap-candidate}.  A high-precision non-interval diagnostic
gave a residual of approximately \(1.88\times10^{-14}\); this diagnostic is
not used in the proof.

Only the candidate is truncated.  The nonlinear residual is not projected
onto the first \(N\) modes.  In particular,
\begin{equation}
 \cos(ku)\sin(\ell u)
 =\frac12\bigl(\sin((k+\ell)u)+\sin((\ell-k)u)\bigr),
 \label{eq:supp-cap-product-modes}
\end{equation}
so products such as \(yy_u\) and \(yH[y_u]\) can contain modes through
degree \(2N=300\).  The implementation forms the complete products.  The
non-polynomial functions \(D\), \(D^{-1}\), and \(G\) are not replaced by
finite Fourier truncations at all; they are represented by interval Taylor
data and global derivative bounds.

\subsection{Path A: exact finite Fourier algebra}
\label{sec:supp-cap-path-a}

For a trigonometric polynomial
\begin{equation}
 F(u)=\frac{\alpha_0}{2}
 +\sum_{k=1}^{K}\bigl(\alpha_k\cos(ku)+\gamma_k\sin(ku)\bigr),
 \label{eq:supp-cap-trigpoly}
\end{equation}
Path A constructs
\begin{equation}
 H[F](u)=\sum_{k=1}^{K}
 \bigl(\alpha_k\sin(ku)-\gamma_k\cos(ku)\bigr).
 \label{eq:supp-cap-path-a-action}
\end{equation}
Differentiation and multiplication are likewise carried out by finite
coefficient formulas.  The inputs are Arb balls, so ``exact'' here means
exact at the level of the finite Fourier representation: every arithmetic
rounding error is enclosed by Arb, and there is no Fourier truncation or
aliasing error attached to the operation.

Applying Eq.~\eqref{eq:supp-cap-path-a-action} to
\begin{equation}
 y_u(u)=-\sum_{k=1}^{N}k a_k\sin(ku)
\end{equation}
gives
\begin{equation}
 H[y_u](u)=\sum_{k=1}^{N}k a_k\cos(ku).
 \label{eq:supp-cap-Hyu}
\end{equation}
This explicit function is denoted \(H_y\) in the implementation.  It is
used directly in \(D\), \(G\), and the local product \(yH_y\).

For any derivative order \(p\geq1\), the same coefficient data give the
rigorous global estimate
\begin{equation}
 \|F^{(p)}\|_\infty
 \leq\sum_{k=1}^{K}k^p\bigl(|\alpha_k|+|\gamma_k|\bigr),
 \label{eq:supp-cap-trig-sup}
\end{equation}
with the constant term included separately when \(p=0\).  These finite sums
are evaluated with outward rounding and seed the derivative-bound chain used
by both pathways.

The corresponding Lean files are
\begin{center}
 \texttt{Definitions.lean},\qquad \texttt{HilbertOnTrigPoly.lean}.
\end{center}
The first defines the Fourier-multiplier operator.  The second proves the
sine and cosine rules, compatibility with finite sums, and the bridge theorem
\begin{center}
 \texttt{periodicHilbert\_yu\_eq\_Hyu}.
\end{center}
Thus the explicit finite expression in Eq.~\eqref{eq:supp-cap-Hyu} is proved
to equal the abstract Fourier-defined Hilbert transform of \(y_u\).

\subsection{Path B: certified discrete Hilbert transform}
\label{sec:supp-cap-path-b}

Let \(F\) be a smooth real \(2\pi\)-periodic function and let
\(u_j=2\pi j/N_{\rm grid}\), where \(N_{\rm grid}\) is even.  The discrete
Hilbert transform has Fourier multipliers \(-i\) on positive discrete
frequencies, \(+i\) on negative frequencies, and zero at the constant and
Nyquist modes.  Julia realizes this operator as the circular convolution
\begin{equation}
 H_{N_{\rm grid}}[F](u_j)
 =\sum_{n=0}^{N_{\rm grid}-1}h_n
 F(u_{j-n\bmod N_{\rm grid}}),
 \label{eq:supp-cap-cot-convolution}
\end{equation}
with
\begin{equation}
 h_n=
 \begin{cases}
 \displaystyle\frac{2}{N_{\rm grid}}
 \cot\!\left(\frac{\pi n}{N_{\rm grid}}\right),&n\text{ odd},\\[1.2ex]
 0,&n\text{ even}.
 \end{cases}
 \label{eq:supp-cap-cot-kernel}
\end{equation}
The cotangent formula is not a second definition of the continuous Hilbert
transform.  It is the inverse discrete Fourier transform of the discrete
multiplier, and therefore an efficient real-space implementation of the
same finite-dimensional operator.

The circular convolution is computed as a linear polynomial convolution
followed by folding modulo \(x^{N_{\rm grid}}-1\).  Arb complex-ball
polynomial multiplication encloses all arithmetic and transform rounding.
The use of a fast convolution algorithm affects runtime, not the operator or
the proof.

The sampled transform still differs from the continuous Hilbert transform
because Fourier modes separated by a multiple of \(N_{\rm grid}\) have the
same samples.  Write
\begin{equation}
 \widehat F_k=\frac{1}{2\pi}\int_0^{2\pi}F(u)e^{-iku}\,du.
\end{equation}
If \(F\) is periodic and
\(\|F^{(p)}\|_\infty\leq C_p\) for an integer \(p\geq2\), integration by
parts \(p\) times gives
\begin{equation}
 |\widehat F_k|\leq\frac{C_p}{|k|^p},\qquad k\ne0.
 \label{eq:supp-cap-fourier-decay}
\end{equation}
The discrete Fourier coefficient at a principal frequency is the sum of all
continuous coefficients aliased into that class.  The continuous and
discrete Hilbert multipliers differ by at most two on an aliased mode.
Combining this observation with Eq.~\eqref{eq:supp-cap-fourier-decay}, the
two-sided Fourier tail, and the integral test yields
\begin{equation}
 \bigl|H[F](u_j)-H_{N_{\rm grid}}[F](u_j)\bigr|
 \leq
 \varepsilon_{\rm alias}(F)
 :=\frac{4C_p}
 {(p-1)(N_{\rm grid}/2-1)^{p-1}}.
 \label{eq:supp-cap-aliasing}
\end{equation}
The symmetric interval
\([ -\varepsilon_{\rm alias},\varepsilon_{\rm alias} ]\) is added to every
output ball.  The resulting balls enclose the continuous Hilbert transform
at the grid nodes, not merely its discrete approximation.

The transform is applied at every derivative level needed by the final
Taylor argument.  Since the periodic Hilbert transform commutes with
differentiation under the smoothness hypotheses used here,
\begin{equation}
 \partial_u^mH[F]=H[\partial_u^mF],
 \label{eq:supp-cap-H-commutes}
\end{equation}
the same calculation supplies local enclosures of all required derivatives.
Global derivative bounds for the output use
\begin{equation}
 \|H[F]\|_\infty\leq\frac{\pi^2}{3}\|F''\|_\infty,
 \label{eq:supp-cap-H-tail}
\end{equation}
which follows by two integrations by parts and
\(2\sum_{k\geq1}k^{-2}=\pi^2/3\).

The Lean file \texttt{AliasingErrorBound.lean} proves the DFT sampling
identity, Fourier-decay and tail estimates, the bound
Eq.~\eqref{eq:supp-cap-aliasing}, and the equivalence between the multiplier
form and the cotangent convolution in Eq.~\eqref{eq:supp-cap-cot-convolution}.
The file \texttt{FourierTailHilbert.lean} proves
Eq.~\eqref{eq:supp-cap-H-tail} and the supporting Fourier/Hilbert facts.
Finally, \texttt{FFTHilbertCorrectness.lean} packages these results into
correctness theorems for one transform and for the derivative-indexed Taylor
data returned by Julia.

\subsection{Interval Taylor data}
\label{sec:supp-cap-taylor-data}

An Arb ball \([m\mathbin{\pm}r]\) represents the closed interval
\([m-r,m+r]\).  Elementary operations such as addition, multiplication,
division away from zero, square root on a positive interval, sine, and cosine
are rounded outward.  Repeated use of the interval inclusion principle then
guarantees that the final ball contains the exact result for the supplied
dyadic inputs.

The rigor of the full verifier is not automatic from pointwise ball
arithmetic alone.  The code also supplies mathematical bounds for variation
between grid points, unresolved Fourier modes, and nonlinear derivative
growth.  These analytic corrections are computed as Arb balls and inserted
explicitly into the enclosure.

For a function \(F\), Julia's \texttt{TaylorSeries} object stores
\begin{equation}
 F^{(m)}(u_j)\in I_{j,m}(F),
 \quad 0\leq m\leq M,
 \qquad
 C_m(F)\geq\|F^{(m)}\|_\infty,
 \quad 0\leq m\leq M+1.
 \label{eq:supp-cap-taylor-record}
\end{equation}
Products use the exact Leibniz identity
\begin{equation}
 (FQ)^{(m)}
 =\sum_{r=0}^{m}\binom mr F^{(r)}Q^{(m-r)},
 \label{eq:supp-cap-leibniz}
\end{equation}
both for local derivative balls and for global bounds.  Reciprocals and
square roots are propagated through Fa\`a di Bruno formulas.  A downward
grid-and-Lipschitz pass can tighten an existing bound by
\begin{equation}
 \|F^{(m)}\|_\infty
 \leq
 \max_j\sup\{|x|:x\in I_{j,m}(F)\}
 +hC_{m+1}(F),
 \qquad h=\frac{2\pi}{N_{\rm grid}}.
 \label{eq:supp-cap-tightening}
\end{equation}
The new bound is retained only when its outward upper endpoint improves the
previous rigorous estimate.

\subsection{\texorpdfstring{Worked evaluation of \(yH[y_u]\)}
{Worked evaluation of yHyu}}
\label{sec:supp-cap-yhyu-worked}

This subsection spells out one representative nonlinear term in the
residual, namely
\begin{equation}
 P(u)=y(u)H[y_u](u).
 \label{eq:supp-cap-P-yhyu}
\end{equation}
The notation is multiplicative: \(P\) is the pointwise product of \(y\) and
the Hilbert transform of \(y_u\).  It is not the term \(H[yy_u]\), which is
computed separately.  This distinction matters because \(yH[y_u]\) is a
local nonlinear product, while \(H[yy_u]\) applies a nonlocal Fourier
multiplier after multiplication.

For the certified candidate,
\begin{equation}
 y(u)=\frac{a_0}{2}+\sum_{k=1}^{N}a_k\cos(ku),
 \qquad N=\CandidateDegree.
\end{equation}
Path A differentiates and applies the periodic Hilbert transform
mode-by-mode, giving
\begin{equation}
 y_u(u)=-\sum_{k=1}^{N}ka_k\sin(ku),
 \qquad
 H[y_u](u)=\sum_{k=1}^{N}ka_k\cos(ku).
 \label{eq:supp-cap-yu-Hyu-worked}
\end{equation}
Thus \(P\) is itself a finite even trigonometric polynomial.  Formally,
\begin{equation}
 P(u)=
 \left(\frac{a_0}{2}+\sum_{k=1}^{N}a_k\cos(ku)\right)
 \left(\sum_{\ell=1}^{N}\ell a_\ell\cos(\ell u)\right),
 \label{eq:supp-cap-yhyu-product-expanded}
\end{equation}
and the identity
\begin{equation}
 \cos(ku)\cos(\ell u)
 =\frac12\bigl(\cos((k+\ell)u)+\cos((k-\ell)u)\bigr)
\end{equation}
shows that modes through degree \(2N\) may occur.  No projection back to
degree \(N\) is made.

The production verifier evaluates Eq.~\eqref{eq:supp-cap-P-yhyu} through
interval Taylor data.  Let
\[
 u_j=\frac{2\pi(j-1)}{N_{\rm grid}},
 \qquad
 j=1,\ldots,N_{\rm grid},
 \qquad
 N_{\rm grid}=\VerificationGrid.
\]
For \(F=y\) and \(F=H[y_u]\), Julia first constructs a Taylor record
consisting of Arb balls
\begin{equation}
 F^{(m)}(u_j)\in I_{j,m}(F),
 \qquad 0\leq m\leq M,
 \label{eq:supp-cap-yhyu-input-records}
\end{equation}
and global derivative bounds
\begin{equation}
 C_m(F)\geq\|F^{(m)}\|_\infty,
 \qquad 0\leq m\leq M+1.
 \label{eq:supp-cap-yhyu-input-supbounds}
\end{equation}
For the certificate \(M=\TaylorOrder\).  The values in
Eq.~\eqref{eq:supp-cap-yhyu-input-records} are not ordinary floating-point
samples: each entry is obtained by evaluating a finite trigonometric
polynomial, or one of its coefficient derivatives, using Arb enclosures of
the coefficients and outward-rounded sine and cosine operations.  The bounds in
Eq.~\eqref{eq:supp-cap-yhyu-input-supbounds} come from the finite
coefficient estimate in Eq.~\eqref{eq:supp-cap-trig-sup}.

The Taylor record for \(P=yH[y_u]\) is then built by the Leibniz rule.  For
each grid point and each \(0\leq m\leq M\),
\begin{equation}
 P^{(m)}(u_j)\in
 I_{j,m}(P):=
 \sum_{r=0}^{m}\binom mr
 I_{j,r}(y)\,I_{j,m-r}(H[y_u]),
 \label{eq:supp-cap-yhyu-local-leibniz}
\end{equation}
where the products and sums on the right are interval operations rounded
outward.  The corresponding initial global derivative bounds satisfy
\begin{equation}
 \|P^{(m)}\|_\infty\leq C_m^{\rm seed}(P):=
 \sum_{r=0}^{m}\binom mr
 C_r(y)\,C_{m-r}(H[y_u]),
 \qquad 0\leq m\leq M+1.
 \label{eq:supp-cap-yhyu-sup-leibniz}
\end{equation}
The sum is evaluated outward, and its upper endpoint supplies the stored
bound.  After this product record is formed, the grid-and-Lipschitz pass
computes
\begin{equation}
 G_m(P)=\max_j\sup\{|x|:x\in I_{j,m}(P)\},\qquad 0\leq m\leq M,
\end{equation}
and walks downwards from \(m=M\) to \(m=0\), implementing the update
\begin{equation}
 C_{M+1}(P)=C_{M+1}^{\rm seed}(P),\qquad
 C_m(P)=\min\{C_m^{\rm seed}(P),\,G_m(P)+hC_{m+1}(P)\}.
 \label{eq:supp-cap-yhyu-tightening}
\end{equation}
Here the minimum denotes comparison of outward upper endpoints, and
\(C_{m+1}(P)\) is the already updated next-order bound.  The mean value
theorem on each cell justifies the second candidate.  This pass changes
only the global bound vector; the local derivative balls remain unchanged.

For the production order \(M=6\), the pass can improve orders \(6,5,\ldots,0\)
but leaves \(C_7(P)\) unchanged.  This distinction determines its role in
the final certificate: \(P\) enters the residual additively, and the final
order-six scan uses its local derivative balls and its order-seven bound.
Consequently, tightening the lower bounds of this particular product does
not improve the final residual estimate in the current pipeline.  The same
tightening routine can affect the final estimate elsewhere, when tightened
bounds feed subsequent products, nonlinear compositions, or truncation of a
higher-order Taylor record.

Only after the product Taylor record has been assembled does the verifier
use Taylor's theorem on a cell.  For \(u=u_j+t\), \(0\leq t\leq h\), with
\(h=2\pi/N_{\rm grid}\), the product alone satisfies
\begin{equation}
 |P(u)|\leq
 \sum_{m=0}^{M}
 \frac{U_{j,m}(P)h^m}{m!}
 +\frac{C_{M+1}(P)h^{M+1}}{(M+1)!},
 \label{eq:supp-cap-yhyu-cell-bound}
\end{equation}
where \(U_{j,m}(P)=\sup\{|x|:x\in I_{j,m}(P)\}\).  In the residual
calculation, \(P\) is not scanned in isolation for the final theorem; its
row-wise derivative balls are subtracted into the residual Taylor record,
and the final scan is applied to the assembled residual.  Equation
\eqref{eq:supp-cap-yhyu-cell-bound} is nevertheless the local continuum
principle that justifies the product record itself.

\paragraph{Julia pipeline for this term.}
The monolithic route in \texttt{ResidualBound.jl} constructs the term as
follows:
\begin{verbatim}
y   = cosine_series(y_coeffs; prec=prec)
yu  = deriv(y)
Hyu = hilbert(yu)

ts_y_h = trigpoly_to_taylor(y, N_grid, M_h; prec=prec)
ts_Hyu = trigpoly_to_taylor(Hyu, N_grid, M_order; prec=prec)
ts_y   = truncate_order(ts_y_h, M_order)

ts_y_Hyu = ts_y * ts_Hyu
ts_R = Arb(mu; prec=prec) * ts_Hyu - ts_y - ts_y_Hyu - ts_Hyyu
\end{verbatim}
The first three lines are finite Fourier algebra.  The two
\texttt{trigpoly\_to\_taylor} calls evaluate derivative towers at the global
grid points and attach derivative sup-norm bounds.  The multiplication line
is the Leibniz construction in
Eq.~\eqref{eq:supp-cap-yhyu-local-leibniz}; it does not multiply two
already-scanned cell supremum bounds.  Internally, multiplication ends with
\texttt{return tighten\_sup\_bounds!(out)} in \texttt{TaylorSeries.jl}.

The sharded Crunchy2 route performs the same operations on contiguous
blocks of rows.  The following excerpt uses the setup record \texttt{su},
precision \texttt{p}, and accumulator defined by that driver:
\begin{verbatim}
for si in shard_range
    jr = su.shards[si]
    ts_y_h = trigpoly_to_taylor(su.y, su.N_grid, su.M_h;
                               prec=p, j_range=jr)
    ts_Hyu = trigpoly_to_taylor(su.Hyu, su.N_grid, M_order;
                               prec=p, j_range=jr)
    ts_y   = truncate_order(ts_y_h, M_order)

    ts_y_Hyu = ts_y * ts_Hyu
    _acc!(acc_yHyu, ts_y_Hyu.derivs, M_order)
    write_matrix_slices(store, "yHyu", si, ts_y_Hyu.derivs)
end
\end{verbatim}
The accumulator records per-derivative maxima needed for the later
grid-and-Lipschitz tightening.  In the scalar replay stage, the same
Leibniz formula is applied to the global sup-bound vectors, then tightened
using maxima merged over the entire grid:
\begin{verbatim}
sb_y_Hyu = _leibniz_sup_bounds(sb_y, sb_Hyu,
                              M_order, p, su.binom_M)
_tighten_from_max!(sb_y_Hyu, st.acc_yHyu.ubs,
                   M_order, su.h_arb, p)
\end{verbatim}
These are the global bounds used in the replay; maxima from a single shard
alone do not certify a bound on the entire period.  The loop inside
\texttt{\_tighten\_from\_max!} visits \texttt{M:-1:0}, leaving
\texttt{sb\_y\_Hyu[M+2]} unchanged.  In the final scan, the stored \texttt{yHyu} rows
are read back and subtracted in the same algebraic order as the monolithic
residual assembly.

\paragraph{Lean interface for this term.}
Lean defines the functions appearing above as finite sums:
\begin{align}
 y(a,u)&=\sum_k a_k\cos(ku),\\
 y_u(a,u)&=-\sum_k k a_k\sin(ku),\\
 H_y(a,u)&=\sum_k k a_k\cos(ku).
\end{align}
The theorem \texttt{periodicHilbert\_yu\_eq\_Hyu} identifies the explicit
finite sum \(H_y\) with the abstract Fourier-defined periodic Hilbert
transform of \(y_u\).  Lean's sum uses the full constant coefficient; the
generated half-DC certificate therefore stores \(a_0/2\) in its zeroth
coefficient to match Julia's convention.

The generic proof in \texttt{CertificateAxioms.lean} asks its assumed Arb
interface for Taylor records of \(y\) and \(H_y\), and applies the
Taylor-series multiplication certificate.  Schematically (with angle
brackets denoting Lean tuple destructuring and omitted arguments denoted by
ellipses):
\begin{verbatim}
obtain <ts_y, h_y, hh_y> :=
  cert_arb_trigpoly_to_taylor (y a_cert) M N_fft hN_fft 0 ...

obtain <ts_Hyu, h_Hyu, hh_Hyu> :=
  cert_arb_trigpoly_to_taylor (Hyu a_cert) M N_fft hN_fft 0 ...

obtain <ts_yHyu, h_yHyu, hh_yHyu> :=
  cert_arb_ts_mul h_y h_Hyu (hh_y.trans hh_Hyu.symm)
\end{verbatim}
The resulting semantic statement is that \texttt{ts\_yHyu} is a valid
Taylor record for
\[
 u\mapsto y(a_{\rm cert},u)\,H_y(a_{\rm cert},u).
\]
During residual assembly this generic proof subtracts the record after forming
\(\mu H_y-y\).  Thus the formal proof tracks the same operator order as the
Julia computation: \(y\), \(H_y\), product \(yH_y\), then subtraction into
the residual.

The dedicated half-DC theorem for the two-million-point run uses a different
numerical interface, in the module
\texttt{PaperCertificateAxioms}:
it assumes the two numerical certificates
\texttt{cert\_arb\_dmin\_bound} and
\texttt{cert\_arb\_sup\_bound\_le\_eps0}, then proves periodic extension
of the residual bound in Lean.  It does not independently replay the Arb
product or its tightening pass.  The generic excerpt above describes the
term's Taylor-record interface, not the numerical dependency chain of that
dedicated theorem.

\paragraph{Performance role.}
The archived Crunchy2 logs do not isolate a separate timer for
\texttt{yHyu}.  Its cost is included in the first sharded pointwise stage,
where each worker also computes \(yy_u\), denominator data, reciprocal data,
and the capillary auxiliary records.  For each grid row, the
\texttt{yHyu} product uses the \(M+1\) stored derivative columns and the
triangular Leibniz sums, hence \(O(M^2)\) Arb operations per grid point
with \(M=\TaylorOrder\).  At order six this is 28 Leibniz summands per row,
or 56 million across the full grid.  The grid reduction for tightening
visits seven derivative columns, or 14 million entries; its scalar downward
pass has seven steps.  These operation counts do not determine separate
wall times, which were not recorded for this term.

In the production run the eight first-stage workers ran concurrently.  Their
wall times ranged from 1:56:25 to 2:14:58, with peak resident memory between
4.972 GB and 9.324 GB.  The subsequent serial stages took 1:22.08 for the
merge, 21:33.92 for the certified FFT/Hilbert stage, and 11:49.78 for the
final residual scan.  Since the first-stage workers overlap, the effective
post-warmup wall time is approximately the slowest worker plus these serial
stages, about 2:49:44.  Sharding keeps \texttt{yHyu} as 1024-row Taylor
blocks in memory during stage 1 and as exact Arb slices on disk afterward;
the full two-million-row product matrix is never required to reside in RAM
at once.

\subsection{Positivity of the capillary denominator}
\label{sec:supp-cap-denominator}

The reciprocal terms are not evaluated until \(D\) has been proved
positive.  Set
\begin{equation}
 A=y_u,\qquad B=1+H[y_u],\qquad D=\sqrt{A^2+B^2}.
\end{equation}
At each grid point, Arb computes a lower endpoint \(\underline D_j\) for
\(D(u_j)\).  The reverse triangle inequality gives
\begin{equation}
 |D(u)-D(v)|
 \leq |A(u)-A(v)|+|B(u)-B(v)|.
\end{equation}
Consequently, if
\(L_A\geq\|A'\|_\infty\) and \(L_B\geq\|B'\|_\infty\), then
\begin{equation}
 d_{\rm Lip}:=
 \min_j\bigl(\underline D_j-h(L_A+L_B)\bigr)
 \leq\inf_uD(u).
 \label{eq:supp-cap-D-lipschitz}
\end{equation}
This proof neither differentiates nor divides by \(D\), so it does not
presuppose the positivity it establishes.  The computation stops unless
outward-rounded arithmetic proves \(d_{\rm Lip}>0\).

After positivity is known, Julia constructs Taylor data for \(D\) and
obtains a second lower bound from the local derivatives and the Taylor
remainder.  The implementation selects the ball having the larger rigorous
lower endpoint.  Both candidates are valid lower bounds; selecting the
larger one only sharpens the subsequent reciprocal estimates.

\subsection{\texorpdfstring{Higher derivatives of \(D\) and \(G\)}
{Higher derivatives of D and G}}
\label{sec:supp-cap-DG-derivatives}

Direct Fa\`a di Bruno estimates for many derivatives of a square root and
reciprocal can become unnecessarily large.  The verifier therefore also
uses algebraic recurrences.  Let
\(S=A^2+B^2=D^2\).  Differentiating \(D^2=S\) \(m\) times and isolating the
two endpoint terms gives
\begin{equation}
 2D D^{(m)}
 =S^{(m)}-
 \sum_{r=1}^{m-1}\binom mrD^{(r)}D^{(m-r)}.
 \label{eq:supp-cap-D-recurrence}
\end{equation}
If \(D\geq d>0\), then
\begin{equation}
 \|D^{(m)}\|_\infty
 \leq\frac{1}{2d}\left(
 \|S^{(m)}\|_\infty+
 \sum_{r=1}^{m-1}\binom mr
 \|D^{(r)}\|_\infty\|D^{(m-r)}\|_\infty
 \right).
 \label{eq:supp-cap-D-bound}
\end{equation}
Similarly, the identity \(GD=B\) gives
\begin{equation}
 G^{(m)}D
 =B^{(m)}-
 \sum_{r=0}^{m-1}\binom mrG^{(r)}D^{(m-r)},
 \label{eq:supp-cap-G-recurrence}
\end{equation}
and therefore
\begin{equation}
 \|G^{(m)}\|_\infty
 \leq\frac{1}{d}\left(
 \|B^{(m)}\|_\infty+
 \sum_{r=0}^{m-1}\binom mr
 \|G^{(r)}\|_\infty\|D^{(m-r)}\|_\infty
 \right).
 \label{eq:supp-cap-G-bound}
\end{equation}
The quantities \(S\) and \(B\) are finite trigonometric polynomials, so
their derivative bounds come directly from Eq.~\eqref{eq:supp-cap-trig-sup}.
The recurrence avoids replacing \(D\) or \(G\) by a finite Fourier series.
Julia compares these algebraic bounds with the valid composition-based
bounds and uses the tighter outward endpoint where needed for the top
derivatives in Path B.

\subsection{Residual assembly and passage from the grid to the torus}
\label{sec:supp-cap-grid-to-continuum}

The local interval records are composed in the same order as
Eq.~\eqref{eq:supp-cap-residual}: finite data for \(y\), \(y_u\), and
\(H[y_u]\); Leibniz products; the Path-B transform of \(yy_u\); construction
of \(D\), \(D^{-1}\), and \(G\); differentiation of \(y_u/D\) and \(G\);
the Path-B transform of \(\partial_uG\); and the final linear combination.
For every grid node and every \(0\leq m\leq M\), the result is an interval
\begin{equation}
 R^{(m)}(u_j)\in I_{j,m}(R),
 \label{eq:supp-cap-R-grid}
\end{equation}
together with
\(C_{M+1}(R)\geq\|R^{(M+1)}\|_\infty\).

For \(u=u_j+t\), \(0\leq t\leq h\), Taylor's theorem yields
\begin{equation}
 |R(u)|\leq
 \sum_{m=0}^{M}
 \frac{U_{j,m}(R)h^m}{m!}
 +\frac{C_{M+1}(R)h^{M+1}}{(M+1)!},
 \label{eq:supp-cap-cell-bound}
\end{equation}
where
\(U_{j,m}(R)=\sup\{|x|:x\in I_{j,m}(R)\}\).  The outward-rounded maximum of
the right-hand side over all cells is the certified \(\varepsilon_*\).  This
is the step that converts finitely many computations into a statement about
every point of the continuous periodic domain.  With \(M=\TaylorOrder\),
the final remainder scales as \(h^7\).

\subsection{Julia/Arb implementation}
\label{sec:supp-cap-julia}

The package \texttt{BabenkoRigorous} implements the numerical proof.  Its
principal files have the following roles.
\begin{description}
 \item[\texttt{FourierSeries.jl}.]
 Defines the finite \texttt{TrigPoly} representation and the Path-A
 operations: evaluation, differentiation, Hilbert transform,
 multiplication, and coefficient-based supremum bounds.

 \item[\texttt{TaylorSeries.jl}.]
 Defines local derivative matrices and global derivative-bound vectors;
 implements Taylor supremum and infimum scans, Leibniz products,
 reciprocal and square-root composition, downward bound tightening,
 denominator positivity, and the algebraic \(D/G\) recurrences.

 \item[\texttt{GDerivatives.jl}.]
 Constructs the Taylor representation of
 \(G=(1+H[y_u])/D\) from the certified denominator and reciprocal.

 \item[\texttt{FFTHilbert.jl}.]
 Implements Path B by Arb complex-polynomial convolution with the cotangent
 kernel, adds Eq.~\eqref{eq:supp-cap-aliasing} to each output ball, and lifts
 the transform through all derivative levels using
 Eq.~\eqref{eq:supp-cap-H-commutes} and the global bound
 Eq.~\eqref{eq:supp-cap-H-tail}.

 \item[\texttt{ResidualBound.jl}.]
 Assembles the complete gravity--capillary residual and applies the final
 Taylor scan.  In the implemented route, \(H[yy_u]\) and
 \(H[\partial_uG]\) both pass through \texttt{fft\_hilbert\_series}; only the
 latter intrinsically requires the non-polynomial pathway.

 \item[\texttt{ShardStore.jl} and \texttt{Sharded.jl}.]
 Reorganize the same computation into exact on-disk shards so that a large
 grid does not require all Taylor matrices to be resident simultaneously.

 \item[\texttt{Certificate.jl}.]
 Converts binary inputs and outward endpoints to exact rational literals and
 emits the generated Lean data module.  The current formatter emits no
 axioms; the two numerical anchors live in a separate certificate-specific
 Lean module.
\end{description}

\paragraph{Verifier order.}
The production verifier follows the order below.  This order is useful for
audit because every later division or Hilbert transform consumes bounds
created earlier in the list.
\begin{enumerate}
 \item Read the \CandidateCount{} stored coefficients and construct
 \texttt{TrigPoly} objects for \(y\), \(y_u\), and \(H[y_u]\).
 \item Convert the finite series into Taylor records on the
 \(N_{\rm grid}\)-point grid, retaining enough derivative orders to support
 the final \(M=\TaylorOrder\) Taylor scan.
 \item Form the polynomial products \(yH[y_u]\) and \(yy_u\) by the Leibniz
 rule, and route \(yy_u\) through the certified sampled Hilbert transform.
 \item Prove a positive lower bound for
 \(D=\sqrt{y_u^2+(1+H[y_u])^2}\) before any reciprocal is evaluated.
 \item Build \(D^{-1}\), \(G=(1+H[y_u])/D\), \(\partial_u(y_u/D)\), and
 \(\partial_uG\) as Taylor records with outward-rounded derivative bounds.
 \item Apply the certified sampled Hilbert transform to \(\partial_uG\),
 using the supplied high-derivative bounds to add the aliasing radius and
 the Fourier-tail derivative bounds.
 \item Assemble the residual record and scan all cells with
 Eq.~\eqref{eq:supp-cap-cell-bound}; write the residual ball, the
 denominator lower ball, and the generated certificate.
\end{enumerate}

Most local arithmetic operations are ordinary point operations lifted to
Arb balls, but the proof also depends on the explicitly programmed analytic
corrections in Eqs.~\eqref{eq:supp-cap-aliasing},
\eqref{eq:supp-cap-tightening}, \eqref{eq:supp-cap-D-lipschitz}, and
\eqref{eq:supp-cap-cell-bound}.  Arb automatically encloses numerical
rounding; it does not invent these mathematical bounds.

\subsection{Memory-sharded computation}
\label{sec:supp-cap-sharding}

At \(N_{\rm grid}=\VerificationGrid\), storing every intermediate matrix at
once is unnecessarily expensive.  The sharded implementation exploits the
structure of the calculation:
\begin{enumerate}
 \item Compute coefficient-level derivative bounds once from the finite
 trigonometric data.
 \item Pointwise nonlinear recurrences and min/max reductions are streamed
 over contiguous grid shards.
 \item Each global Hilbert transform is processed one derivative column at a
 time.  Exact Arb column slices are stored on disk.
 \item The final residual is reassembled shard by shard and the Taylor-cell
 maximum is reduced in ascending grid order.
\end{enumerate}
Intermediate balls are serialized with Arb's exact dump/load representation,
including mantissa, exponent, and radius; decimal text is not used for this
round trip.  The reduction order matches the monolithic loops, and repository
tests compare the sharded and monolithic paths on smaller grids.  Sharding
therefore changes memory use and scheduling, not the mathematical enclosure.

The successful two-million-point calculation used shards of 1024 points.
The first stage ran on eight workers with four Julia threads per worker; the
main transform and final stages used eight threads.  The stages consisted of
parallel pointwise reductions, merging of the worker states,
derivative-by-derivative certified convolutions, and a final sharded residual
scan.  Separate operating-system processes between major stages released the
memory associated with completed stages and left exact state dumps for
reproducibility.

\subsection{Lean definitions and reusable analytic library}
\label{sec:supp-cap-lean-library}

The Lean project \texttt{WaveCAP} defines the mathematical operator rather
than treating it as an uninterpreted symbol.  For a real function \(f\),
\texttt{Definitions.lean} sets
\begin{equation}
 H[f](u)=\operatorname{Re}\sum_{n\in\mathbb Z}
 \bigl(-i\,\operatorname{sgn}(n)\bigr)\widehat f_n e^{inu},
 \label{eq:supp-cap-lean-H}
\end{equation}
with multiplier zero at \(n=0\).  The same file defines the finite profile,
its derivative, the explicit \(H_y\), \(D\), \(G\), and \(R\) from
Eq.~\eqref{eq:supp-cap-residual}.  The coefficient type
\(\operatorname{Fin}(N)\to\mathbb R\) makes every profile sum finite.

The reusable formal library mirrors the proof architecture:
\begin{itemize}
\raggedright
 \item \textbf{Hilbert transform.}
 \texttt{HilbertOnTrigPoly.lean} proves Path A and
 \(H[y_u]=H_y\).  \texttt{FourierTailHilbert.lean} proves the
 \(\pi^2/3\) estimate and mode facts.
 \texttt{HilbertCommutesDeriv.lean} proves derivative commutation.

 \item \textbf{Aliasing and Path B.}
 \texttt{AliasingErrorBound.lean} proves the sampled DFT aliasing theorem and
 the cot-kernel equivalence.  The file
 \texttt{FFTHilbert}\allowbreak\texttt{Correctness.lean} states the
 correctness contract for the interval transform and its Taylor-series lift.

 \item \textbf{Derivative calculus.}
 \texttt{LeibnizBound.lean} and \texttt{FaaDiBrunoBound.lean} prove the
 product and composition bounds used by Julia.

 \item \textbf{Taylor data.}
 \texttt{TaylorRemainder.lean} and
 \texttt{GridToSupNormTaylor.lean} prove the local remainder and the
 grid-to-continuum theorem.  The file
 \texttt{TaylorSeries}\allowbreak\texttt{Correctness.lean} defines
 the semantic invariant \texttt{IsTaylorSeriesOf} and proves that the
 arithmetic operations preserve it.  \texttt{SeriesCorrectness.lean} covers
 the square-root and reciprocal recurrences, while
 \texttt{TightenSupBounds}\allowbreak\texttt{Lemma.lean} justifies the
 downward tightening pass.

 \item \textbf{Denominator and capillary chain.}
 \texttt{DminBound}\allowbreak\texttt{Correctness.lean} formalizes the
 positivity calculation, and \texttt{TightDGSupBounds.lean} proves
 Eqs.~\eqref{eq:supp-cap-D-bound}--\eqref{eq:supp-cap-G-bound}.

 \item \textbf{Residual assembly.}
 \texttt{ResidualAssembly}\allowbreak\texttt{Correctness.lean} composes the
 Path-A and Path-B records into the full residual, proves periodicity, and
 proves the general extension of a one-period bound to every real argument.
\end{itemize}
These are generic theorems about the mathematics and data-structure
contracts of the pipeline.  They are useful independently of the particular
numerical profile.

\subsection{Exact certificate data and convention alignment}
\label{sec:supp-cap-exact-data}

The generated module
\begin{center}
\texttt{CertificatePhase4ScalarN150}\allowbreak
\texttt{Phase420260827HalfDC.lean}
\end{center}
contains \CandidateCount{} exact rational coefficients and exact rational
representations of \(\mu\), \(\beta\), \(N_{\rm grid}\), \(M\), the
residual upper endpoint, and the denominator lower endpoint.  In particular,
\begin{align}
 \mu_{\rm cert}
 &=\frac{1538048895284245}{1125899906842624},\\
 \beta_{\rm cert}
 &=\frac{2233821443972785}{9007199254740992},
 \label{eq:supp-cap-exact-parameters}
\end{align}
and
\begin{align}
 \varepsilon_{0,\rm cert}
 &=\frac{66570650447636094343693612891855023841}
 {374144419156711147060143317175368453031918731001856},
 \label{eq:supp-cap-exact-eps}\\
 D_{\rm lower,cert}
 &=\frac{30024018999569627600479976468787845960585}
 {87112285931760246646623899502532662132736}.
 \label{eq:supp-cap-exact-D}
\end{align}
Thus moving the certificate into Lean introduces no decimal-rounding step.

Julia and Lean use different storage conventions for the constant cosine
mode.  Julia evaluates the first source entry as \(a_0/2\), while the generic
Lean function \texttt{y} uses its zeroth stored entry directly.  The
certificate therefore stores
\begin{equation}
 \widehat a_0=a_0/2,
 \qquad
 \widehat a_k=a_k\quad(k\geq1).
 \label{eq:supp-cap-half-DC}
\end{equation}
The generic module \texttt{HalfDCConvention.lean} proves the two coefficient
cases and proves that this adapter preserves \(y_u\), \(H_y\), \(D\), and
\(G\).  The certificate-specific theorem checks the exact zeroth-entry
identity.  The source vector is not also encoded as a second complete Lean
vector, so the present formal theorem does not compare two named
151-component vectors; the generated certificate already contains the
converted vector.

The audit module \texttt{PaperAudit.lean} separately proves
the sign identity in Eq.~\eqref{eq:supp-cap-sign} and the corresponding
equality of absolute values.  Both the half-DC identities and the residual
sign bridge are independent of the numerical trust anchors.

\subsection{The numerical trust boundary in Lean}
\label{sec:supp-cap-trust-boundary}

The present certificate does not serialize the millions of local derivative
balls, all global derivative vectors, or the intermediate convolution
enclosures needed for Lean to reconstruct the two-million-point run.  The
certificate-specific numerical interface therefore consists of exactly two
assumptions recording the final Arb conclusions:
\begin{align}
 \mathsf A_D:\quad&
 D_{\rm lower,cert}>0
 \quad\text{and}\quad
 \forall u\in\mathbb R,\quad
 D_{\rm lower,cert}\leq D(a_{\rm cert};u),
 \label{eq:supp-cap-anchor-D}\\
 \mathsf A_R:\quad&
 \mathsf A_D\Longrightarrow
 \forall u\in[0,2\pi],\quad
 |R(a_{\rm cert},\mu_{\rm cert},\beta_{\rm cert};u)|
 \leq\varepsilon_{0,\rm cert}.
 \label{eq:supp-cap-anchor-R}
\end{align}
Their Lean names are
\begin{center}
 \texttt{cert\_arb\_dmin\_bound},\qquad
 \texttt{cert\_arb\_sup\_bound\_le\_eps0}.
\end{center}
They are declared in
\begin{center}
 \texttt{CertificatePhase4ScalarN150}\allowbreak
 \texttt{Phase420260827HalfDCAxioms.lean}.
\end{center}

This distinction is important.  Lean proves the universal analytic
ingredients used by the verifier, but the current certificate-facing theorem
does not instantiate those generic results from serialized grid matrices.
Instead, it imports Eqs.~\eqref{eq:supp-cap-anchor-D}--
\eqref{eq:supp-cap-anchor-R} as the explicit numerical trust boundary.
Accordingly, the result should be described as a machine-checked theorem
with two externally computed Arb anchors, not as a replay of every interval
operation inside Lean.

From the two anchors Lean performs the remaining deductions:
\begin{enumerate}
 \item It projects the global denominator lower bound from \(\mathsf A_D\).
 \item It applies \(\mathsf A_R\) on \([0,2\pi]\).
 \item It proves periodicity of the residual and extends the bound to all
 \(u\in\mathbb R\).
 \item It verifies by exact rational normalization that
 \(\varepsilon_{0,\rm cert}<178/10^{15}\) and
 \(3446588/10^7<D_{\rm lower,cert}\).
 \item It transfers the absolute residual theorem to the manuscript operator
 using Eq.~\eqref{eq:supp-cap-sign}.
\end{enumerate}
The corresponding public modules are
\begin{center}
 \texttt{MainPhase4ScalarN150}\allowbreak
 \texttt{Phase420260827HalfDC.lean},\\
 \texttt{PaperAudit.lean}.
\end{center}

The axiom audit for the exact and decimal residual theorems reports the
standard Mathlib principles \texttt{propext}, \texttt{Classical.choice}, and
\texttt{Quot.sound}, together with precisely the two numerical anchors above.
The denominator theorem uses only the first numerical anchor.  The half-DC
and sign-alignment results use no certificate-specific numerical axioms.

\subsection{Numerical result and interpretation of precision}
\label{sec:supp-cap-result}

The successful Crunchy2 computation used
\begin{align}
 N&=\CandidateDegree,
 &N_{\rm grid}&=\VerificationGrid,\\
 M&=\TaylorOrder,
 &p_{\rm Arb}&=\WorkingPrecision\ \text{bits},
 \label{eq:supp-cap-run-parameters}
\end{align}
with shard size 1024.  The final Arb ball was reported as
\begin{align}
 \operatorname{mid}(B_R)&=\ResidualMidpoint,\\
 \operatorname{rad}(B_R)&=\ResidualRadius,\\
 \inf B_R&=\ResidualLower,\\
 \sup B_R&=\ResidualUpper.
 \label{eq:supp-cap-final-ball}
\end{align}
Here \(B_R\) is the ball returned by the supremum-bound computation; it is
not an enclosure of the range of \(R\), which explains why its lower endpoint
may be negative.  The theorem uses the outward upper endpoint
\(\ResidualUpper\), not the midpoint or radius separately.  The denominator
calculation returned the outward lower endpoint
\begin{equation}
 D_{\rm lower}=\DLower.
\end{equation}
The run log ends with \texttt{ALL\_STAGES\_COMPLETE}, and the final stage has
zero exit status.

The phrase ``near machine precision'' refers here to the practical accuracy
floor of the double-precision candidate and the complete nonlinear residual,
not to the rounding unit of the 128-bit Arb calculation.  IEEE double
precision has machine epsilon \(2^{-52}\approx2.22\times10^{-16}\) and unit
roundoff \(2^{-53}\approx1.11\times10^{-16}\).  Cancellation, derivative
amplification, conditioning, and the storage of the candidate coefficients
can raise the attainable residual above these raw constants.  The 128-bit
Arb precision, roughly 38 decimal digits, makes arithmetic rounding in the
verification negligible relative to the analytic Taylor and aliasing
corrections; it does not force the final certificate radius to be
\(10^{-38}\).

\newpage
\subsection{Reproducibility and file map}
\label{sec:supp-cap-reproducibility}

The paths and production commands in this section preserve the original
research checkout and Crunchy2 environment. In the paper's publication
repository, \texttt{BabenkoRigorous/} corresponds to \texttt{Julia/},
\texttt{WaveCAP/WaveCAP/} corresponds to \texttt{Lean/WaveCAP/}, and
\texttt{wave-cap/WaveCAP/} corresponds to \texttt{Lean/}. The repository
README gives build and verification commands for the publication layout.

\begin{samepage}
The archived run directory
\begin{center}
\texttt{BabenkoRigorous/results/phase4\_scalar\_n150\_20260827\_2m\_20260829}
\end{center}
contains the run description and checksums for the raw logs, stage logs,
coefficient file, Arb endpoint dumps, scripts, and generated Lean
certificate.  The installed half-DC certificate was compared byte-for-byte
with the generated artifact.  The dedicated Lean theorem and audit modules
were checked directly with the pinned Lean and Mathlib version
\texttt{v4.29.0-rc6}.
\end{samepage}

The principal archived artifacts have the following SHA-256 checksums:
\begin{verbatim}
4d65357a64a0ccfbf0ed67f8dd0e1970f0ceb549a8cbf647d8490e4ac48abb07
ee8fe5db1ea3bbd6e37d0888fdadc00c9c9e1d6ce81cb9d46ae3e4db8ffcdf43
\end{verbatim}
for the run archive and the half-DC alignment archive, respectively.  The
installed generated half-DC Lean certificate has SHA-256 checksum
\begin{verbatim}
2190a52f3520ab0b15683eab592ee1105797d737a22f2e24a869130f5a431ed2
\end{verbatim}

\paragraph{Crunchy2 execution table.}
The following table is transcribed from the archived
\texttt{/usr/bin/time -v} logs and from the stage summaries printed by
\texttt{stage\_run\_phase4\_scalar\_n150\_20260827.jl}.  The eight
\texttt{p1part} workers were launched concurrently, so their wall times
should not be summed.  The total production wall time after warmup is
approximately the slowest \texttt{p1part} worker plus the three serial
stages.

\begingroup
\setlength{\tabcolsep}{3pt}
\begin{longtable}{@{}>{\raggedright\arraybackslash}p{0.25\linewidth}>{\raggedright\arraybackslash}p{0.11\linewidth}>{\raggedright\arraybackslash}p{0.15\linewidth}>{\raggedright\arraybackslash}p{0.15\linewidth}>{\raggedright\arraybackslash}p{0.13\linewidth}@{}}
\caption{Crunchy2 stage timing and memory for the certified two-million-point run.}
\label{tab:supp-crunchy2-runtime}\\
\toprule
Log file & Threads & Wall time & Julia elapsed & Peak RSS \\
\midrule
\endfirsthead
\toprule
Log file & Threads & Wall time & Julia elapsed & Peak RSS \\
\midrule
\endhead
\bottomrule
\endfoot
\path{p1part_1of8.log} & 4 & 1:57:25 & 6959.9 s & 5.019 GB \\
\path{p1part_2of8.log} & 4 & 2:14:58 & 8007.8 s & 9.242 GB \\
\path{p1part_3of8.log} & 4 & 1:56:25 & 6895.6 s & 4.972 GB \\
\path{p1part_4of8.log} & 4 & 2:07:07 & 7541.7 s & 9.287 GB \\
\path{p1part_5of8.log} & 4 & 2:13:04 & 7897.6 s & 9.324 GB \\
\path{p1part_6of8.log} & 4 & 1:58:34 & 7032.1 s & 5.088 GB \\
\path{p1part_7of8.log} & 4 & 2:09:00 & 7655.0 s & 9.267 GB \\
\path{p1part_8of8.log} & 4 & 2:13:52 & 7944.0 s & 9.283 GB \\
\path{p1merge.log} & 8 & 1:22.08 & 0.7 s & 0.481 GB \\
\path{fft.log} & 8 & 21:33.92 & 1212.4 s & 5.149 GB \\
\path{s3.log} & 8 & 11:49.78 & 628.2 s & 0.477 GB \\
\end{longtable}
\endgroup

\newpage
The final \texttt{s3} log reports
\begin{verbatim}
bound midpoint = +1.9928311263803702e-14
bound radius   = 1.5799937398052128e-13
bound interval = [-1.3807106271671758e-13, +1.7792768524432498e-13]
bound ubound   = 1.7792768524432498e-13
D_lower lbound = 3.4465883518530394e-01
\end{verbatim}
and the top-level run log ends with \texttt{ALL\_STAGES\_COMPLETE}.

\paragraph{Production command.}
On \texttt{crunchy2}, from the \texttt{BabenkoRigorous} directory, the
archived run is reproduced by
\begin{verbatim}
WORK=/scratch/ms14479/phase4_scalar_n150_20260827_2m_retry_work
LOGDIR=/scratch/ms14479/phase4_scalar_n150_20260827_2m_retry_logs
RUNTIME=apptainer
SIF=/home/tb97/transit-cap-ccf/cap-ccf.sif
BINDPATH=/scratch:/scratch
WORKER_THREADS=4
MAIN_THREADS=8
scripts/p1_parallel_run_phase4_scalar_n150_20260827.sh \
  "$WORK" 1024 2000000 8 "$LOGDIR"
\end{verbatim}
The first run writes the non-half-DC generated certificate and the two Arb
dumps \path{staged_bound.arbdump} and
\path{staged_dlower.arbdump}.  The half-DC Lean certificate used by the
paper is then re-emitted from the existing dumps, without rerunning the
two-million-point interval calculation, by
\begin{verbatim}
apptainer exec --bind /scratch:/scratch \
  /home/tb97/transit-cap-ccf/cap-ccf.sif \
  julia --startup-file=no --compiled-modules=no --project=. \
  scripts/reemit_phase4_n150_half_dc_certificate.jl "$WORK"
\end{verbatim}

\paragraph{Lean verification commands.}
After the generated half-DC certificate is installed in
\path{WaveCAP/WaveCAP/}, the certificate-facing targets are checked by
\begin{verbatim}
cd wave-cap/WaveCAP
lake build
lake env lean WaveCAP/PaperTheorems.lean
lake env lean WaveCAP/PaperAudit.lean
\end{verbatim}
The direct \texttt{lake env lean} calls are included because the dedicated
half-DC headline and audit modules are not exercised by the default library
root in the current repository layout.

For readers following the implementation, the Julia source names below are
relative to \texttt{BabenkoRigorous/}.  The Lean source names are relative
to \texttt{WaveCAP/WaveCAP/}.
\begin{itemize}
\raggedright
 \item \textbf{Path A.}
 Julia: \texttt{src/FourierSeries.jl}.
 Lean: \texttt{Definitions.lean} and \texttt{HilbertOnTrigPoly.lean}.

 \item \textbf{Path B and the cot kernel.}
 Julia: \texttt{src/FFTHilbert.jl}.
 Lean: \texttt{AliasingErrorBound.lean},
 \texttt{FourierTailHilbert.lean}, and
 \texttt{FFTHilbertCorrectness.lean}.

 \item \textbf{Nonlinear Taylor data and denominator.}
 Julia: \texttt{src/TaylorSeries.jl} and
 \texttt{src/GDerivatives.jl}.
 Lean: \texttt{TaylorSeriesCorrectness.lean},
 \texttt{SeriesCorrectness.lean}, \texttt{DminBoundCorrectness.lean}, and
 \texttt{TightDGSupBounds.lean}.

 \item \textbf{Residual and continuum bound.}
 Julia: \texttt{src/ResidualBound.jl}.
 Lean: \texttt{TaylorRemainder.lean},
 \texttt{GridToSupNormTaylor.lean}, and
 \texttt{ResidualAssemblyCorrectness.lean}.

 \item \textbf{Large-grid organization.}
 Julia: \texttt{src/ShardStore.jl}, \texttt{src/Sharded.jl}, and
 \texttt{scripts/stage\_run\_phase4\_scalar\_n150\_20260827.jl}.

 \item \textbf{Certificate and theorem.}
 Julia: \texttt{src/Certificate.jl}.
 Lean: the generated half-DC certificate, its certificate-specific axiom
 module, the corresponding \texttt{Main} module, and
 \texttt{PaperAudit.lean}.
\end{itemize}

These components yield a reproducible continuum residual bound for the
profile in Eq.~\eqref{eq:supp-cap-candidate}.  They state the
mathematical pathways, numerical enclosures, formal definitions, and
remaining numerical trust boundary separately.
\endgroup

\end{document}
